\documentclass[10pt,letterpaper]{article}
\usepackage{technical-testing}
\setDocTitle{Animation Playtest Standard Operating Procedure}
\setDocSubtitle{Controlled execution and review procedure}
\setDocVersion{1.0}
\setDocStatus{Educational / Draft}
\setDocOwner{Game Development}
\setDocDate{2026-10-30}
\setDocSummary{Controlled execution and review procedure}
\begin{document}
\makedoctitle
\makedocsummary
\tableofcontents\clearpage

\section{Research aim}
Determine whether players detect, understand, and can respond to animated cues as intended. A useable study does not require a large formal laboratory, but it does require consent, a consistent script, a controlled build and a plan for handling feedback.
\section{Preparation}
\begin{enumerate}
\item Identify a single decision, such as attack anticipation duration or pickup confirmation latency.
\item Freeze animation clips and game rules except the tested timing parameter. Label A/B/C by timing configuration without steering observers.
\item Recruit representative participants and define exclusions and accessibility accommodations; avoid sensitive biometric data collection unless justified and approved.
\item Prepare a short neutral instruction, task sequence, observable metrics, and an opt-out path.
\item Pilot the test and verify telemetry (input timestamp, cue time, response, error, FPS and state trace).
\end{enumerate}
\section{Execution}
Capture task success and incorrect responses. Observe anticipation, action readability, perceived weight, motion comfort, interface feedback and lost control. Separate participant quotation from researcher inference. Randomize configuration order where learning effects matter. Avoid asking leading questions such as whether a new animation is ``better.''
\section{Interpretation}
Compute median and 95th percentile latency only when sample size and instrumentation justify them; show counts and distributions, not only averages. Segment device, settings and input scheme where material. Report uncertainty, outliers, build constraints and artistic trade-offs.
\section{Exit conditions}
A signed research record identifies tested version, anonymized observations, decision owner, change request, and follow-up verification. Use paired playtest and accessibility scorecards.
\end{document}
